% Source-faithful mathematical blueprint of Hartshorne, Chapter III (Cohomology).
% Authorial exposition, examples, and exercises are omitted; definitions, statements, and proof
% arguments are retained. Statement coordinates are the printed chapter coordinates.
\section{Derived functors}
\label{ha-ch3-sec1}

\begin{definition}[Abelian category, complexes, and exact functors]
\label{ha-ch3-def-abelian}
\source{hartshorne-algebraic-geometry:III.1}
An abelian category has abelian Hom groups, finite direct sums, kernels and
cokernels, and the usual normality and factorization axioms.  A complex is a
sequence $d^{i}:A^{i}\to A^{i+1}$ with $d^{i+1}d^i=0$; its cohomology is
$\ker d^i/\operatorname{im}d^{i-1}$.  A functor is left exact if it preserves
the left three terms of every short exact sequence.  An object is injective if
$\operatorname{Hom}(-,I)$ is exact, and an injective resolution is an exact
embedding into a complex of injectives.  An object is $F$-acyclic if
$R^iF(A)=0$ for $i>0$; a category has enough injectives if every object embeds
in an injective object.  For a left exact functor $F$, the derived functors are
$R^iF(A)=H^i(F(I^\bullet))$ for an injective resolution $I^\bullet$.
\end{definition}

\begin{theorem}[Derived functors]
\label{ha-ch3-thm-derived}
\dcref{III.1.1A}
\source{hartshorne-algebraic-geometry:III.1.1A}
If $\mathcal A$ has enough injectives and $F:\mathcal A\to\mathcal B$ is
additive left exact, then $R^iF(A)=H^i(F(I^\bullet))$ is independent of the
injective resolution and is functorial and additive.  There is a natural
isomorphism $R^0F\cong F$, and each short exact sequence
$0\to A'\to A\to A''\to0$ gives a functorial long exact sequence with
connecting maps $\delta^i:R^iF(A'')\to R^{i+1}F(A')$.  Morphisms of short
exact sequences commute with the connecting maps, and $R^iF(I)=0$ for
$i>0$ and $I$ injective.

\uses{ha-ch3-def-abelian}
\end{theorem}
\begin{proof}

Comparison maps between injective resolutions are unique up to homotopy, so
applying $F$ gives canonically isomorphic cohomology.  The short exact
sequence is replaced by a short exact sequence of resolutions; the long
exact cohomology sequence gives the connecting maps.  Homotopic comparison
maps induce the same cohomology maps.  An injective has the resolution
$0\to I\to I\to0$, proving the vanishing.

\end{proof}

\begin{proposition}[Acyclic resolutions]
\label{ha-ch3-prop-acyclic-resolution}
\dcref{III.1.2A}
\source{hartshorne-algebraic-geometry:III.1.2A}
If $0\to A\to J^0\to J^1\to\cdots$ is exact and every $J^i$ is
$F$-acyclic, then $R^iF(A)\cong H^i(F(J^\bullet))$ naturally.

\uses{ha-ch3-thm-derived,ha-ch3-def-abelian}
\end{proposition}
\begin{proof}

Embed the acyclic resolution and an injective resolution into a common
comparison diagram.  The quotient complexes are acyclic after applying $F$
by induction and the long exact sequence, hence the comparison is a
quasi-isomorphism.

\end{proof}

\begin{definition}[Cohomological $\delta$-functor and universality]
\label{ha-ch3-def-delta}
\source{hartshorne-algebraic-geometry:III.1}
A covariant $\delta$-functor is functors $T^i$ with connecting maps giving
long exact sequences and commuting with morphisms of short exact sequences.
It is universal if a map in degree zero to another $\delta$-functor extends
uniquely in every degree.  A functor is effaceable if every object embeds into
an object on which it vanishes.
\end{definition}

\begin{theorem}[Effaceability criterion]
\label{ha-ch3-thm-effaceable}
\dcref{III.1.3A}
\source{hartshorne-algebraic-geometry:III.1.3A}
An effaceable positive-degree $\delta$-functor is universal.

\uses{ha-ch3-def-delta}
\end{theorem}
\begin{proof}

Embed successively in objects killing the given positive-degree functor and
use the long exact sequence to define the extension degree by degree.  The
connecting maps force uniqueness.

\end{proof}

\begin{corollary}[Universal property of derived functors]
\label{ha-ch3-cor-universal-derived}
\dcref{III.1.4}
\source{hartshorne-algebraic-geometry:III.1.4}
For a left exact $F$ on a category with enough injectives, $(R^iF)$ is a
universal $\delta$-functor with degree zero $F$.  Conversely, every universal
$\delta$-functor has left exact degree zero and is naturally isomorphic to the
derived functors of that degree-zero functor.

\uses{ha-ch3-def-delta,ha-ch3-thm-effaceable,ha-ch3-thm-derived}
\end{corollary}
\begin{proof}

Every object embeds in an injective, so positive derived functors are
effaceable; apply III.1.3A.  The converse follows by applying universality to
the derived functors of $T^0$.

\end{proof}

\section{Cohomology of sheaves}
\label{ha-ch3-sec2}
\begin{proposition}[Enough injectives for modules]
\label{ha-ch3-prop-module-injective}\dcref{III.2.1A}
\source{hartshorne-algebraic-geometry:III.2.1A}
Every module over a ring embeds in an injective module.

\uses{ha-ch3-def-abelian}
\end{proposition}
\begin{proof}

Embed the module in its injective envelope, or equivalently apply Baer's criterion: extend the
identity map from the module to an injective hull and regard the hull as the required injective module.

\end{proof}
\begin{proposition}[Enough injective sheaves]
\label{ha-ch3-prop-sheaf-injective}\dcref{III.2.2}
\source{hartshorne-algebraic-geometry:III.2.2}
For every ringed space $(X,\mathcal O_X)$, the category of $\mathcal O_X$-modules
has enough injectives.

\uses{ha-ch3-def-abelian}
\end{proposition}
\begin{proof}

Embed each stalk $\mathcal F_x$ into an injective module $I_x$ and embed
$\mathcal F$ into the product of the skyscraper direct images $j_{x*}I_x$.
The target is injective because its Hom functor is a product of exact stalk
functors followed by exact Hom functors.

\end{proof}
\begin{corollary}[Enough injective abelian sheaves]
\label{ha-ch3-cor-sheaf-injective}\dcref{III.2.3}
\source{hartshorne-algebraic-geometry:III.2.3}
The category of sheaves of abelian groups on any topological space has enough
injectives.

\uses{ha-ch3-prop-sheaf-injective}
\end{corollary}
\begin{proof}

For an abelian sheaf $\mathcal F$, embed each stalk $\mathcal F_x$ in an injective abelian group
$I_x$.  The map sending a section to its germs embeds $\mathcal F$ in the product of the skyscraper
sheaves $i_{x*}I_x$.  A product of skyscraper injectives is injective: a map from a subsheaf extends
stalkwise, and the sheaf axiom assembles the extensions.

\end{proof}
\begin{definition}[Sheaf cohomology and flasque sheaves]
\label{ha-ch3-def-cohomology-flasque}
\source{hartshorne-algebraic-geometry:III.2}
For an abelian sheaf $\mathcal F$, $H^i(X,\mathcal F)$ is the $i$th right
derived functor of $\Gamma(X,-)$.  A sheaf is flasque when all restriction
maps are surjective.
\uses{ha-ch3-def-abelian,ha-ch3-thm-derived}
\end{definition}
\begin{lemma}[Injectives are flasque]
\label{ha-ch3-lem-injective-flasque}\dcref{III.2.4}
\source{hartshorne-algebraic-geometry:III.2.4}
Every injective $\mathcal O_X$-module is flasque.

\uses{ha-ch3-def-cohomology-flasque}
\end{lemma}
\begin{proof}

For $V\subset U$, the inclusion of extension-by-zero structure sheaves
$\mathcal O_V\hookrightarrow\mathcal O_U$ and injectivity give a surjection
$\operatorname{Hom}(\mathcal O_U,I)\to\operatorname{Hom}(\mathcal O_V,I)$,
which is $I(U)\to I(V)$.

\end{proof}
\begin{proposition}[Flasque acyclicity]
\label{ha-ch3-prop-flasque-acyclic}\dcref{III.2.5}
\source{hartshorne-algebraic-geometry:III.2.5}
If $\mathcal F$ is flasque, then $H^i(X,\mathcal F)=0$ for $i>0$.

\uses{ha-ch3-def-cohomology-flasque,ha-ch3-lem-injective-flasque}
\end{proposition}
\begin{proof}

Embed $\mathcal F$ in an injective $I$ and let $Q$ be the quotient.  Both
$I$ and $Q$ are flasque, and $0\to\Gamma(F)\to\Gamma(I)\to\Gamma(Q)\to0$
is exact.  The long exact sequence and induction on degree give the claim.

\end{proof}
\begin{proposition}[Global sections as derived functors]
\label{ha-ch3-prop-cohomology-derived}\dcref{III.2.6}
\source{hartshorne-algebraic-geometry:III.2.6}
The derived functors of $\Gamma(X,-)$ on $\mathcal O_X$-modules are the
groups $H^i(X,-)$ defined on abelian sheaves.

\uses{ha-ch3-def-cohomology-flasque,ha-ch3-thm-derived,ha-ch3-prop-acyclic-resolution}
\end{proposition}
\begin{proof}
Use an injective resolution, which is flasque by III.2.4 and
acyclic by III.2.5, and apply III.1.2A.
\end{proof}
\begin{theorem}[Grothendieck vanishing]
\label{ha-ch3-thm-grothendieck-vanishing}\dcref{III.2.7}
\source{hartshorne-algebraic-geometry:III.2.7}
If $X$ is noetherian of dimension $n$, then $H^i(X,\mathcal F)=0$ for every
abelian sheaf $\mathcal F$ and $i>n$.

\end{theorem}
\begin{proof}

Induct on dimension and number of irreducible components.  Reduce by the
closed/open exact sequence to an irreducible $X$.  Express a sheaf as a
filtered direct limit of finitely generated subsheaves; cohomology commutes
with these limits because limits of flasques are flasque on a noetherian
space.  Reduce a finitely generated sheaf to a quotient of a constant sheaf
on an open set.  The kernel is treated by noetherian induction on its support;
the constant sheaf on an irreducible space is flasque, and the closed
complement has smaller dimension.

\end{proof}
\begin{lemma}[Limits of flasque sheaves]
\label{ha-ch3-lem-limits-flasque}\dcref{III.2.8}
\source{hartshorne-algebraic-geometry:III.2.8}
On a noetherian space, a filtered direct limit of flasque sheaves is flasque.

\end{lemma}
\begin{proof}

For $V\subseteq U$, choose a finite affine (hence quasi-compact) cover of $U$ refining a cover on
which a representative section is defined.  Surjectivity of each restriction map in the system gives
compatible lifts on this finite cover; the sheaf axiom glues them.  Thus every section of the colimit
over $V$ lifts to $U$.

\end{proof}
\begin{proposition}[Cohomology and direct limits]
\label{ha-ch3-prop-cohomology-direct-limits}\dcref{III.2.9}
\source{hartshorne-algebraic-geometry:III.2.9}
For a filtered direct system of abelian sheaves on a noetherian space,
$H^i(X,\varinjlim\mathcal F_\alpha)\cong\varinjlim H^i(X,\mathcal F_\alpha)$;
in particular cohomology commutes with arbitrary direct sums.

\end{proposition}
\begin{proof}

The canonical maps form morphisms of delta-functors.  Efface both sides by
the functorial sheaf of discontinuous sections; the limit is flasque by
III.2.8, so universality (III.1.3A) identifies the two delta-functors.

\end{proof}
\begin{lemma}[Closed-subset extension]
\label{ha-ch3-lem-closed-extension}\dcref{III.2.10}
\source{hartshorne-algebraic-geometry:III.2.10}
For a closed inclusion $j:Y\hookrightarrow X$,
$H^i(Y,\mathcal F)\cong H^i(X,j_*\mathcal F)$.

\end{lemma}
\begin{proof}
Apply a flasque resolution on $Y$ and extend each term by zero.
\end{proof}

\section{Noetherian affine schemes}
\label{ha-ch3-sec3}
\begin{proposition}[Krull's theorem]
\label{ha-ch3-prop-krull}\dcref{III.3.1A}
\source{hartshorne-algebraic-geometry:III.3.1A}
If $A$ is noetherian, $M\subset N$ are finite $A$-modules and $\mathfrak a$
is an ideal, the $\mathfrak a$-adic topology on $M$ is induced from $N$.

\end{proposition}
\begin{proof}

The book cites the Artin--Rees lemma (Atiyah--Macdonald, or Bourbaki) for
this assertion; no independent proof is supplied here.

\end{proof}
\begin{lemma}[Torsion in an injective module]
\label{ha-ch3-lem-torsion-injective}\dcref{III.3.2}
\source{hartshorne-algebraic-geometry:III.3.2}
For an injective $A$-module $I$, the submodule $\Gamma_{\mathfrak a}(I)$ of
$\mathfrak a$-power torsion is injective.

\end{lemma}
\begin{proof}

For a map from an ideal $\mathfrak b$ to the torsion submodule, finite
generation and Krull's theorem make it factor through
$\mathfrak b/(\mathfrak b\cap\mathfrak a^n)$.  Extend this map to $A/\mathfrak a^n$
using injectivity of $I$; its image is torsion, proving Baer's criterion.

\end{proof}
\begin{lemma}[Localization of injectives]
\label{ha-ch3-lem-localize-injective}\dcref{III.3.3}
\source{hartshorne-algebraic-geometry:III.3.3}
For injective $I$ over a noetherian ring, $I\to I_f$ is surjective for every
$f\in A$.

\end{lemma}
\begin{proof}

The ascending annihilators of powers of $f$ stabilize.  Represent an element
of $I_f$ by $y/f^r$, extend the resulting map from the stabilized principal
ideal to $A$ by injectivity, and the image of $1$ maps to the given element.

\end{proof}
\begin{proposition}[Injective quasi-coherent sheaves are flasque]
\label{ha-ch3-prop-affine-injective-flasque}\dcref{III.3.4}
\source{hartshorne-algebraic-geometry:III.3.4}
If $I$ is injective over noetherian $A$, then $\widetilde I$ on
$\operatorname{Spec}A$ is flasque.

\end{proposition}
\begin{proof}

Use noetherian induction on the support.  Restriction to a principal open is
surjective by III.3.3; the difference is supported on the closed complement,
where its torsion submodule is injective by III.3.2 and induction applies.

\end{proof}
\begin{theorem}[Affine vanishing]
\label{ha-ch3-thm-affine-vanishing}\dcref{III.3.5}
\source{hartshorne-algebraic-geometry:III.3.5}
For $X=\operatorname{Spec}A$ noetherian and $\mathcal F$ quasi-coherent,
$H^i(X,\mathcal F)=0$ for $i>0$ and $\Gamma(X,\mathcal F)$ recovers its
module.
\end{theorem}
\begin{proof}
Resolve the module of global sections by injectives, sheafify,
and use III.3.4 and III.2.5; the resulting global-section complex is exact.
\end{proof}
\begin{corollary}[Flasque quasi-coherent embedding]
\label{ha-ch3-cor-flasque-qcoh}\dcref{III.3.6}
\source{hartshorne-algebraic-geometry:III.3.6}
Every quasi-coherent sheaf on a noetherian scheme embeds in a flasque
quasi-coherent sheaf.
\uses{ha-ch3-prop-affine-injective-flasque}
\end{corollary}
\begin{proof}

Cover $X$ by finitely many affine opens.  On each affine, embed the corresponding module in an
injective module and use Proposition~\ref{ha-ch3-prop-affine-injective-flasque}.  The resulting
local embeddings glue after taking their product; the product is quasi-coherent and flasque, and the
stalkwise maps give the required global embedding.

\end{proof}
\begin{theorem}[Serre's affineness criterion]
\label{ha-ch3-thm-serre-affine}\dcref{III.3.7}
\source{hartshorne-algebraic-geometry:III.3.7}
For a noetherian scheme $X$, the following are equivalent: $X$ is affine;
$H^i(X,\mathcal F)=0$ for every quasi-coherent $\mathcal F$ and $i>0$; and
$H^1(X,\mathcal J)=0$ for every coherent ideal sheaf $\mathcal J$.

\end{theorem}
\begin{proof}

The first implication is III.3.5.  For the converse, use vanishing for ideal
sheaves to find, around each closed point, an affine principal open $X_f$.
Choose finitely many $f_i$ covering $X$.  The kernel of
$\mathcal O_X^{\oplus m}\to\mathcal O_X$, $(a_i)\mapsto\sum f_i a_i$,
has a filtration with coherent ideal quotients, so its $H^1$ vanishes.  Global
sections therefore express $1$ as a combination of the $f_i$, and the affine
criterion applies.

\end{proof}

\section{Cech cohomology}
\label{ha-ch3-sec4}
\begin{definition}[Cech complex]
\label{ha-ch3-def-cech}
\source{hartshorne-algebraic-geometry:III.4}
For an ordered open cover $\mathfrak U=(U_i)$, set
$C^p(\mathfrak U,\mathcal F)=\prod_{i_0<\cdots<i_p}\mathcal F(U_{i_0\cdots i_p})$
with the alternating restriction differential, and define
$\check H^p(\mathfrak U,\mathcal F)=H^p(C^\bullet)$.
\end{definition}
\begin{lemma}[Cech degree zero]
\label{ha-ch3-lem-cech-zero}\dcref{III.4.1}
\source{hartshorne-algebraic-geometry:III.4.1}
The canonical map $H^0(\mathfrak U,\mathcal F)\to\Gamma(X,\mathcal F)$ is an
isomorphism.
\end{lemma}
\begin{proof}
A zero cocycle is a compatible family of sections, hence
glues uniquely by the sheaf axiom.
\end{proof}
\begin{lemma}[Sheafified Cech resolution]
\label{ha-ch3-lem-cech-resolution}\dcref{III.4.2}
\source{hartshorne-algebraic-geometry:III.4.2}
The sheafified Cech complex is a resolution of $\mathcal F$.

\end{lemma}
\begin{proof}

Exactness is stalkwise.  At a point choose one cover member containing it;
insertion of that index gives a contracting homotopy in positive degrees.

\end{proof}
\begin{proposition}[Flasque Cech acyclicity]
\label{ha-ch3-prop-cech-flasque}\dcref{III.4.3}
\source{hartshorne-algebraic-geometry:III.4.3}
If $\mathcal F$ is flasque, then $\check H^p(\mathfrak U,\mathcal F)=0$ for
$p>0$.
\end{proposition}
\begin{proof}
The sheafified resolution consists of flasque sheaves, so it
computes both derived and Cech cohomology, and III.2.5 gives vanishing.
\end{proof}
\begin{lemma}[Cech comparison map]
\label{ha-ch3-lem-cech-comparison-map}\dcref{III.4.4}
\source{hartshorne-algebraic-geometry:III.4.4}
There is a natural map $\check H^p(\mathfrak U,\mathcal F)\to H^p(X,\mathcal F)$,
functorial in $\mathcal F$.

\end{lemma}
\begin{proof}

The sheafified Cech resolution maps to an injective resolution of $\mathcal F$ by the comparison
theorem for resolutions.  Applying $\Gamma(X,-)$ gives a cochain map from the Cech complex to an
injective complex; its induced cohomology map is independent of the chosen comparison up to homotopy.

\end{proof}
\begin{theorem}[Cech comparison for affine covers]
\label{ha-ch3-thm-cech-comparison}\dcref{III.4.5}
\source{hartshorne-algebraic-geometry:III.4.5}
If $X$ is noetherian separated, $\mathfrak U$ is an affine open cover, and
$\mathcal F$ is quasi-coherent, then the comparison map is an isomorphism in
every degree.
\end{theorem}
\begin{proof}
Embed $\mathcal F$ in a flasque quasi-coherent sheaf and
induct on degree using the long exact sequences.  All finite intersections are
affine, so their section sequences are exact by III.3.5; the flasque term has
zero positive Cech cohomology by III.4.3.
\end{proof}

\section{Projective space}
\label{ha-ch3-sec5}
\begin{theorem}[Cohomology of projective space]
\label{ha-ch3-thm-projective-space}\dcref{III.5.1}
\source{hartshorne-algebraic-geometry:III.5.1}
Let $X=\mathbf P_A^r$, $r\ge1$, over noetherian $A$.  Then
$\bigoplus_nH^0(X,\mathcal O_X(n))\cong A[x_0,\ldots,x_r]$;
$H^i(X,\mathcal O_X(n))=0$ for $0<i<r$; $H^r(X,\mathcal O_X(-r-1))\cong A$;
and multiplication gives a perfect pairing
$H^0(\mathcal O(n))\times H^r(\mathcal O(-n-r-1))\to A$.

\end{theorem}
\begin{proof}

Use the standard affine cover and the Cech comparison theorem.  The top
Cech quotient is free on Laurent monomials with every exponent negative;
this gives the top group, its grading, and the pairing with nonnegative
monomials.  Localizing at $x_r$ makes the intermediate cohomology vanish;
the hyperplane exact sequence identifies multiplication by $x_r$ with a
bijective shift in intermediate degrees, while localization already makes
each class $x_r$-power torsion.  Thus those groups vanish.

\end{proof}
\begin{theorem}[Serre finiteness and eventual vanishing]
\label{ha-ch3-thm-serre-finiteness}\dcref{III.5.2}
\source{hartshorne-algebraic-geometry:III.5.2}
For $X$ projective over noetherian $A$ with a very ample $\mathcal O_X(1)$ and
$\mathcal F$ coherent, $H^i(X,\mathcal F)$ is finite over $A$, and
$H^i(X,\mathcal F(n))=0$ for $i>0$ and all $n\gg0$.

\end{theorem}
\begin{proof}

Embed $X$ in projective space and use III.5.1 for sums of twists.  Resolve a
coherent sheaf by a finite sum of twists and induct on the cohomological
degree using the long exact sequence; noetherianity preserves finite
generation and the same induction gives eventual vanishing.

\end{proof}
\begin{proposition}[Cohomological criterion for ampleness]
\label{ha-ch3-prop-ample-criterion}\dcref{III.5.3}
\source{hartshorne-algebraic-geometry:III.5.3}
For an invertible sheaf $\mathcal L$ on a proper noetherian $A$-scheme,
$\mathcal L$ is ample iff for every coherent $\mathcal F$,
$H^i(X,\mathcal F\otimes\mathcal L^n)=0$ for $i>0$ and $n\gg0$.

\end{proposition}
\begin{proof}

The forward implication follows after replacing $\mathcal L$ by a very ample
power and applying III.5.2 to finitely many residue classes of exponents.
Conversely apply vanishing to the ideal of each closed point; global sections
surject onto the skyscraper quotient, hence generate every twist locally by
Nakayama.  A finite cover and a common exponent show all coherent sheaves
are eventually generated, which is the ampleness criterion.

\end{proof}

\section{Ext groups and sheaves}
\label{ha-ch3-sec6}
\begin{definition}[Ext]
\label{ha-ch3-def-ext}
\source{hartshorne-algebraic-geometry:III.6}
$\mathcal H om(\mathcal F,\mathcal G)$ is the sheaf of local homomorphisms,
and $\mathcal E xt^i(\mathcal F,\mathcal G)$ is the $i$th right derived
functor of $\mathcal H om(\mathcal F,-)$.  The groups
$\operatorname{Ext}^i_X(\mathcal F,\mathcal G)$ are the right derived
functors of $\operatorname{Hom}_X(\mathcal F,-)$.
\end{definition}
\begin{lemma}[Restriction of injectives]\label{ha-ch3-lem-ext-restriction}\dcref{III.6.1}
\source{hartshorne-algebraic-geometry:III.6.1}
Injective sheaves remain injective after restriction to an open subset.

\end{lemma}
\begin{proof}

Restriction to an open subset has extension by zero as an exact left adjoint.  The adjunction and
injectivity of the original sheaf therefore extend every map into its restriction.

\end{proof}
\begin{proposition}[Restriction of sheaf Ext]\label{ha-ch3-prop-ext-restriction}\dcref{III.6.2}
\source{hartshorne-algebraic-geometry:III.6.2}
For $U\subset X$, $\mathcal E xt_X^i(\mathcal F,\mathcal G)|_U\cong
\mathcal E xt_U^i(\mathcal F|_U,\mathcal G|_U)$.
\uses{ha-ch3-lem-ext-restriction}
\end{proposition}
\begin{proof}

Restrict an injective resolution of $\mathcal G$ and use Lemma~\ref{ha-ch3-lem-ext-restriction}.
Restriction commutes with sheaf Hom, so the two resulting complexes agree on $U$ and have isomorphic
cohomology sheaves.

\end{proof}
\begin{proposition}[Ext from the structure sheaf]\label{ha-ch3-prop-ext-structure}\dcref{III.6.3}
\source{hartshorne-algebraic-geometry:III.6.3}
$\mathcal E xt^0(\mathcal O_X,\mathcal G)=\mathcal G$, higher sheaf Ext from
$\mathcal O_X$ vanishes, and $\operatorname{Ext}^i(\mathcal O_X,\mathcal G)=H^i(X,\mathcal G)$.

\end{proposition}
\begin{proof}

$\mathcal H om(\mathcal O_X,\mathcal G)=\mathcal G$, and $\mathcal O_X$ is locally free of rank one,
so its higher sheaf Ext vanishes.  Taking global sections of an injective resolution of $\mathcal G$
identifies global Ext with sheaf cohomology.

\end{proof}
\begin{proposition}[Long exact Ext sequences]\label{ha-ch3-prop-ext-les}\dcref{III.6.4}
\source{hartshorne-algebraic-geometry:III.6.4}
A short exact sequence in the first variable gives the usual long exact
sequences for both global Ext and sheaf Ext.

\end{proposition}
\begin{proof}
Apply Hom into an injective resolution; Hom into an injective is
exact in the first variable, and use the long exact cohomology sequence.
\end{proof}
\begin{proposition}[Locally free resolutions]\label{ha-ch3-prop-ext-resolution}\dcref{III.6.5}
\source{hartshorne-algebraic-geometry:III.6.5}
If $\mathcal F$ has a locally free finite-rank resolution $\mathcal L_\bullet$,
then $\mathcal E xt^i(\mathcal F,\mathcal G)=H^i(\mathcal H om(\mathcal L_\bullet,\mathcal G))$.

\end{proposition}
\begin{proof}

Apply $\mathcal H om(-,\mathcal G)$ to the resolution.  Since every $\mathcal L_j$ is locally free,
the resulting complex computes the derived Hom object and its cohomology is the stated sheaf Ext.

\end{proof}
\begin{lemma}[Tensoring injectives]\label{ha-ch3-lem-tensor-injective}\dcref{III.6.6}
\source{hartshorne-algebraic-geometry:III.6.6}
If $\mathcal L$ is locally free of finite rank and $I$ injective, then
$\mathcal L\otimes I$ is injective.

\end{lemma}
\begin{proof}

Locally $\mathcal L$ is finite free, so $\mathcal L\otimes I$ is a finite direct sum of copies of
$I$.  Finite direct sums of injectives are injective, and the assertion is local for sheaves.

\end{proof}
\begin{proposition}[Ext and locally free tensor factors]\label{ha-ch3-prop-ext-tensor}\dcref{III.6.7}
\source{hartshorne-algebraic-geometry:III.6.7}
$\operatorname{Ext}^i(\mathcal F\otimes\mathcal L,\mathcal G)\cong
\operatorname{Ext}^i(\mathcal F,\mathcal L^\vee\otimes\mathcal G)$, and likewise for sheaf Ext.

\end{proposition}
\begin{proof}

The adjunction $\mathcal H om(\mathcal F\otimes\mathcal L,\mathcal G)
\cong\mathcal H om(\mathcal F,\mathcal L^\vee\otimes\mathcal G)$ is termwise valid because
$\mathcal L$ is finite locally free.  Resolve $\mathcal G$ injectively and use the preceding lemma.

\end{proof}
\begin{proposition}[Stalks of sheaf Ext]\label{ha-ch3-prop-ext-stalk}\dcref{III.6.8}
\source{hartshorne-algebraic-geometry:III.6.8}
For coherent $\mathcal F$ on a noetherian scheme,
$(\mathcal E xt^i(\mathcal F,\mathcal G))_x\cong
\operatorname{Ext}^i_{\mathcal O_{X,x}}(\mathcal F_x,\mathcal G_x)$.

\end{proposition}
\begin{proof}
Work affine and use a free resolution of the coherent module;
stalks are exact and commute with Hom from finite free modules.
\end{proof}
\begin{proposition}[Ext twisting comparison]\label{ha-ch3-prop-ext-twist}\dcref{III.6.9}
\source{hartshorne-algebraic-geometry:III.6.9}
For coherent $\mathcal F,\mathcal G$ on projective $X$, and fixed $i$, for
$n\gg0$ the natural maps from global sections of $\mathcal E xt^i(\mathcal F,\mathcal G(n))$
to $\operatorname{Ext}^i(\mathcal F,\mathcal G(n))$ are isomorphisms.

\end{proposition}
\begin{proof}
Resolve $\mathcal F$ by locally free sheaves and induct on its
length, using III.5.2 and exactness of sufficiently twisted global sections.
\end{proof}
\begin{proposition}[Projective dimension]\label{ha-ch3-prop-projective-dimension}\dcref{III.6.10A}
\source{hartshorne-algebraic-geometry:III.6.10A}
For a module, projective dimension is the length of its shortest projective
resolution.

\end{proposition}
\begin{proof}

A projective resolution can be shortened exactly when the corresponding syzygy is projective.  Taking
the infimum of the lengths of all projective resolutions gives the invariant, with value $\infty$
when no finite resolution exists.

\end{proof}
\begin{proposition}[Regular local rings]\label{ha-ch3-prop-regular-local}\dcref{III.6.11A}
\source{hartshorne-algebraic-geometry:III.6.11A}
Over a regular local ring, finite projective dimension equals dimension minus
depth.

\end{proposition}
\begin{proof}

A regular local ring has finite global dimension $\dim A$.  The Auslander--Buchsbaum formula gives
$\operatorname{pd}_A M+\operatorname{depth}_A M=\operatorname{depth}A=\dim A$ for every finite module
$M$ of finite projective dimension.

\end{proof}
\begin{proposition}[Ext detects depth]\label{ha-ch3-prop-ext-depth}\dcref{III.6.12A}
\source{hartshorne-algebraic-geometry:III.6.12A}
The Ext groups against the residue field detect the depth of a finite module.

\end{proposition}
\begin{proof}

For $(A,\mathfrak m,k)$ local, take a minimal free resolution of $M$; all differentials have entries
in $\mathfrak m$, so after applying $\operatorname{Hom}_A(-,k)$ the differentials vanish.  The first
nonzero group $\operatorname{Ext}^i_A(k,M)$ is therefore in degree $\operatorname{depth}_A M$.

\end{proof}

\section{Serre duality}
\label{ha-ch3-sec7}
\begin{theorem}[Duality for projective space]\label{ha-ch3-thm-duality-projective}\dcref{III.7.1}
\source{hartshorne-algebraic-geometry:III.7.1}
For $X=\mathbf P_k^n$ and $\omega_X=\mathcal O_X(-n-1)$,
$H^n(X,\omega_X)\cong k$, the pairing
$\operatorname{Hom}(\mathcal F,\omega_X)\times H^n(X,\mathcal F)\to k$
is perfect, and $\operatorname{Ext}^i(\mathcal F,\omega_X)\cong
H^{n-i}(X,\mathcal F)^\vee$ for coherent $\mathcal F$.

\end{theorem}
\begin{proof}
The first assertion and the line-bundle case follow from III.5.1;
resolve a coherent sheaf by sums of twists and use the five lemma for degree
zero and universality of coeffaceable contravariant delta-functors in higher
degrees.
\end{proof}
\begin{definition}[Dualizing sheaf]\label{ha-ch3-def-dualizing}
\source{hartshorne-algebraic-geometry:III.7}
For proper $n$-dimensional $X/k$, a dualizing sheaf is coherent $\omega_X$ with
a trace $H^n(X,\omega_X)\to k$ inducing perfect Hom--top-cohomology pairings.
\end{definition}
\begin{proposition}[Uniqueness of dualizing sheaves]\label{ha-ch3-prop-dualizing-unique}\dcref{III.7.2}
\source{hartshorne-algebraic-geometry:III.7.2}
A dualizing sheaf with trace is unique up to a unique trace-preserving
isomorphism.
\end{proposition}
\begin{proof}
Use the representing property on the two sheaves to get
mutually inverse maps; uniqueness follows from the same property.
\end{proof}
\begin{lemma}[Ext vanishing by codimension]\label{ha-ch3-lem-ext-codim}\dcref{III.7.3}
\source{hartshorne-algebraic-geometry:III.7.3}
If $X\subset\mathbf P^N$ has codimension $r$, then
$\mathcal E xt^i(\mathcal O_X,\omega_{\mathbf P^N})=0$ for $i<r$.

\end{lemma}
\begin{proof}

At a point of $X$, a free resolution of the local defining ideal has length at least the codimension
by the depth inequality.  Applying $\operatorname{Hom}(-,\omega_{\mathbf P^N})$ therefore has no
cohomology below degree $r$; stalkwise vanishing gives the sheaf assertion.

\end{proof}
\begin{lemma}[The dualizing Ext sheaf]\label{ha-ch3-lem-dualizing-ext}\dcref{III.7.4}
\source{hartshorne-algebraic-geometry:III.7.4}
With $\omega_X=\mathcal E xt^r(\mathcal O_X,\omega_{\mathbf P^N})$,
$\operatorname{Hom}_X(\mathcal F,\omega_X)\cong
\operatorname{Ext}_{\mathbf P^N}^r(\mathcal F,\omega_{\mathbf P^N})$.
\uses{ha-ch3-lem-ext-codim}
\end{lemma}
\begin{proof}

Use the local-to-global Ext spectral sequence and the vanishing below $r$ from Lemma
~\ref{ha-ch3-lem-ext-codim}.  Its degree-$r$ edge map is the displayed isomorphism, compatible
with restriction to affine opens and hence global.

\end{proof}
\begin{proposition}[Existence of dualizing sheaves]\label{ha-ch3-prop-dualizing-existence}\dcref{III.7.5}
\source{hartshorne-algebraic-geometry:III.7.5}
Every projective scheme over a field has a dualizing sheaf.

\end{proposition}
\begin{proof}
Embed in projective space and use the Ext sheaf of III.7.4;
projective-space duality supplies the trace and representing property.
\end{proof}
\begin{theorem}[Duality for a projective scheme]\label{ha-ch3-thm-duality-projective-scheme}\dcref{III.7.6}
\source{hartshorne-algebraic-geometry:III.7.6}
For projective equidimensional $n$-fold $X/k$ with dualizing sheaf $\omega_X$,
there are natural maps $\operatorname{Ext}^i(\mathcal F,\omega_X)\to
H^{n-i}(X,\mathcal F)^\vee$.  They are all isomorphisms iff $X$ is
Cohen--Macaulay and equidimensional, equivalently iff every locally free
$\mathcal F$ has $H^i(X,\mathcal F(-q))=0$ for $i<n$ and $q\gg0$.

\end{theorem}
\begin{proof}
Universality gives the maps.  Embed in $\mathbf P^N$ and compare
depth/projective dimension with Ext vanishing to prove the equivalence of the
first two conditions.  Under the vanishing both contravariant delta-functors
are universal, so the maps are isomorphisms; the converse follows by taking
locally free sheaves and applying projective-space duality.
\end{proof}
\begin{corollary}[Cohen--Macaulay duality]\label{ha-ch3-cor-duality-cm}\dcref{III.7.7}
\source{hartshorne-algebraic-geometry:III.7.7}
For Cohen--Macaulay equidimensional projective $X$, duality holds for locally
free sheaves.
\uses{ha-ch3-thm-duality-projective-scheme}
\end{corollary}
\begin{proof}

For a Cohen--Macaulay equidimensional scheme, the depth criterion gives the vanishing hypothesis in
Theorem~\ref{ha-ch3-thm-duality-projective-scheme}; specialize that theorem to locally free
$\mathcal F$.

\end{proof}
\begin{corollary}[Normal vanishing]\label{ha-ch3-cor-normal-vanishing}\dcref{III.7.8}
\source{hartshorne-algebraic-geometry:III.7.8}
If $X$ is normal of dimension at least two, then
$H^i(X,\mathcal F(-q))=0$ for $i<2$ and $q\gg0$.
\uses{ha-ch3-thm-duality-projective-scheme}
\end{corollary}
\begin{proof}

Normality implies Serre's condition $S_2$, so local depth is at least two away from codimension one.
The Ext/depth criterion in Theorem~\ref{ha-ch3-thm-duality-projective-scheme} gives the stated
vanishing after a sufficiently negative twist.

\end{proof}
\begin{corollary}[Connected ample divisors]\label{ha-ch3-cor-connected-divisor}\dcref{III.7.9}
\source{hartshorne-algebraic-geometry:III.7.9}
An ample effective divisor on a normal projective variety of dimension at least
two is connected.

\end{corollary}
\begin{proof}

Apply the normal vanishing statement to the structure sheaf and use the exact
sequence of the divisor; the resulting $H^0$ calculation forces connectedness.

\end{proof}
\begin{proposition}[Koszul regular sequence]\label{ha-ch3-prop-koszul}\dcref{III.7.10A}
\source{hartshorne-algebraic-geometry:III.7.10A}
If $f_1,\ldots,f_r$ is an $M$-regular sequence, the positive homology of its
Koszul complex with coefficients in $M$ vanishes and degree zero is
$M/(f_1,\ldots,f_r)M$.

\end{proposition}
\begin{proof}

The book cites Matsumura (Theorem 43, p. 135) and Serre (IV.A) for the
Koszul exactness assertion; no proof is supplied here.

\end{proof}
\begin{theorem}[Dualizing sheaf of a local complete intersection]\label{ha-ch3-thm-lci-dualizing}\dcref{III.7.11}
\source{hartshorne-algebraic-geometry:III.7.11}
If $X\subset\mathbf P^N$ is a codimension-$r$ local complete intersection with
ideal $\mathcal I$, then
$\omega_X\cong\omega_{\mathbf P^N}\otimes\bigwedge^r(\mathcal I/\mathcal I^2)^\vee$;
it is invertible.
\end{theorem}
\begin{proof}
Locally resolve $\mathcal O_X$ by the Koszul complex of a
regular sequence.  Its dual computes the top Ext and changes by the
determinant under a change of generators, exactly compensated by
$\bigwedge^r(\mathcal I/\mathcal I^2)^\vee$.
\end{proof}
\begin{corollary}[Canonical sheaf on a nonsingular variety]\label{ha-ch3-cor-nonsingular-duality}\dcref{III.7.12}
\source{hartshorne-algebraic-geometry:III.7.12}
For a nonsingular projective variety, $\omega_X$ is the canonical sheaf.
\uses{ha-ch3-thm-lci-dualizing}
\end{corollary}
\begin{proof}

The conormal sequence for a nonsingular closed embedding is locally split with a regular sequence of
equations.  The local-complete-intersection formula of Theorem~\ref{ha-ch3-thm-lci-dualizing}
therefore identifies the dualizing sheaf with $\bigwedge^{\dim X}\Omega_X$.

\end{proof}
\begin{corollary}[Serre duality for nonsingular varieties]\label{ha-ch3-cor-serre-nonsingular}\dcref{III.7.13}
\source{hartshorne-algebraic-geometry:III.7.13}
Serre duality holds with the canonical sheaf; in particular
$H^n(X,\omega_X)\cong k$ and for curves $p_a=p_g$.
\uses{ha-ch3-thm-duality-projective-scheme}
\end{corollary}
\begin{proof}

Apply Theorem~\ref{ha-ch3-thm-duality-projective-scheme} with $\mathcal F=\mathcal O_X$ and use
the preceding identification of $\omega_X$.  For a curve, the equality of the dimensions of
$H^1(X,\mathcal O_X)$ and regular differentials is exactly $p_a=p_g$.

\end{proof}
\begin{theorem}[Existence of residues]\label{ha-ch3-thm-residues-existence}\dcref{III.7.14.1}
\source{hartshorne-algebraic-geometry:III.7.14.1}
Let $X$ be a complete nonsingular curve over an algebraically closed field
$k$, let $K$ be its function field, and let $\Omega_K$ be the space of
meromorphic differentials.  For each closed point $P$ there is a unique
$k$-linear map $\operatorname{res}_P:\Omega_K\to k$ such that, for a local
parameter $t$, $\operatorname{res}_P(t^n\,dt)=0$ for $n\ne-1$ and
$\operatorname{res}_P(t^{-1}dt)=1$.

\end{theorem}
\begin{proof}

The existence and parameter-independence are cited in the text to Serre,
*Algebraic Groups and Class Fields*, Chap. II, and to Tate's construction;
the book does not give a proof here.

\end{proof}
\begin{theorem}[Residue theorem]\label{ha-ch3-thm-residue-theorem}\dcref{III.7.14.2}
\source{hartshorne-algebraic-geometry:III.7.14.2}
For every $\omega\in\Omega_K$, $\sum_{P\in X}\operatorname{res}_P(\omega)=0$.

\end{theorem}
\begin{proof}

The residue theorem is cited in the text to Serre and Tate; no proof is
provided in Hartshorne at this point.

\end{proof}

\section{Higher direct images}
\label{ha-ch3-sec8}
\begin{definition}[Higher direct image]\label{ha-ch3-def-higher-direct-image}
\source{hartshorne-algebraic-geometry:III.8}
$R^if_*\mathcal F$ is the $i$th derived functor of direct image.
\end{definition}
\begin{proposition}[Stalk description]\label{ha-ch3-prop-direct-image-stalk}\dcref{III.8.1}
\source{hartshorne-algebraic-geometry:III.8.1}
The sheaf $R^if_*\mathcal F$ is the sheaf associated to
$U\mapsto H^i(f^{-1}U,\mathcal F)$.

\end{proposition}
\begin{proof}

Take an injective resolution $\mathcal F\to I^\bullet$.  Direct image is left exact, and on an
open $U\subseteq Y$ its cochain complex is $\Gamma(f^{-1}U,I^\bullet)$.  Its cohomology is
$H^i(f^{-1}U,\mathcal F)$, and sheafifying this presheaf gives $R^if_*\mathcal F$.

\end{proof}
\begin{corollary}[Restriction of direct images]\label{ha-ch3-cor-direct-image}\dcref{III.8.2}
\source{hartshorne-algebraic-geometry:III.8.2}
Restriction to an open $V\subset Y$ identifies $(R^if_*\mathcal F)|_V$ with
$R^i(f|_{f^{-1}V})_*\mathcal F$.

\end{corollary}
\begin{proof}

Restrict the complex in the preceding proof.  The sections on an open $W\subseteq V$ are still
$H^i(f^{-1}W,\mathcal F)$, so the associated sheaves agree.

\end{proof}
\begin{corollary}[Flasque vanishing]\label{ha-ch3-cor-flasque-vanishing}\dcref{III.8.3}
\source{hartshorne-algebraic-geometry:III.8.3}
$R^if_*$ vanishes on flasque sheaves.

\end{corollary}
\begin{proof}

For a flasque $\mathcal F$, every inverse-image open $f^{-1}U$ is flasque by restriction.  Its
positive cohomology vanishes, so the stalk description gives $R^if_*\mathcal F=0$ for $i>0$.

\end{proof}
\begin{proposition}[Direct image on module categories]\label{ha-ch3-prop-direct-image-composition}\dcref{III.8.4}
\source{hartshorne-algebraic-geometry:III.8.4}
Let $f:X\to Y$ be a morphism of ringed spaces.  The functors $R^if_*$ on
sheaves of modules are the derived functors of
$f_*:\operatorname{Mod}(X)\to\operatorname{Mod}(Y)$.
\uses{ha-ch3-prop-cohomology-derived}
\end{proposition}
\begin{proof}

This is the derived-functor comparison stated in the book; it follows by
calculating with injective resolutions and Proposition~III.2.6.

\end{proof}
\begin{proposition}[Affine direct images]\label{ha-ch3-prop-affine-direct-image}\dcref{III.8.5}
\source{hartshorne-algebraic-geometry:III.8.5}
Let $X$ be noetherian, let $f:X\to Y$ be a morphism of schemes, and let
$Y=\operatorname{Spec}A$ be affine.  For every quasi-coherent $\mathcal F$ on
$X$ there is a natural isomorphism
\[
R^if_*\mathcal F\cong\widetilde{H^i(X,\mathcal F)}.
\]

\end{proposition}
\begin{proof}

The text proves this by identifying both sides as universal delta-functors,
using the affine case and the vanishing of cohomology for flasque sheaves.

\end{proof}
\begin{corollary}[Quasi-coherence of direct images]\label{ha-ch3-cor-direct-image-affine}\dcref{III.8.6}
\source{hartshorne-algebraic-geometry:III.8.6}
Let $f:X\to Y$ be a morphism of schemes with $X$ noetherian.  If
$\mathcal F$ is quasi-coherent, then each $R^if_*\mathcal F$ is
quasi-coherent.
\uses{ha-ch3-prop-affine-direct-image}
\end{corollary}
\begin{proof}

This follows locally on affine opens of $Y$ from Proposition~III.8.5 and
the quasi-coherent affine calculation.

\end{proof}
\begin{corollary}[Separated affine computation]\label{ha-ch3-cor-separated-direct-image}\dcref{III.8.7}
\source{hartshorne-algebraic-geometry:III.8.7}
In the separated noetherian case, higher direct images are computed on affine
covers by Cech complexes.

\end{corollary}
\begin{proof}

Apply the affine-cover Cech comparison theorem to $f^{-1}U$ for each affine $U\subseteq Y$.
Separatedness makes finite intersections affine, and the identifications are compatible with
restriction in $U$, giving the claimed sheaf calculation.

\end{proof}
\begin{theorem}[Coherence of proper direct images]\label{ha-ch3-thm-direct-image-coherent}\dcref{III.8.8}
\source{hartshorne-algebraic-geometry:III.8.8}
Let $f:X\to Y$ be projective between noetherian schemes, let $\mathcal F$ be
coherent, and let $\mathcal L=\mathcal O_X(1)$ be very ample over $Y$.  Then
(a) for every $n>0$, the natural map
$f^*f_*\mathcal F(n)\to\mathcal F(n)$ is surjective; (b) each
$R^if_*\mathcal F$ is coherent; and (c) $R^if_*\mathcal F(n)=0$ for
$i>0$ and $n\gg0$.

\end{theorem}
\begin{proof}
Reduce to projective space, use the finite generation and eventual
vanishing of III.5.2 on affine opens, and glue the resulting finite modules.
\end{proof}

\section{Flat morphisms}
\label{ha-ch3-sec9}
\begin{proposition}[Flatness criteria]\label{ha-ch3-prop-flat-criteria}\dcref{III.9.1A}
\source{hartshorne-algebraic-geometry:III.9.1A}
For an $A$-module $M$ the following hold: (a) $M$ is flat iff
$-\otimes_A M$ is exact; (b) base extension preserves flatness; (c) flatness
is transitive; (d) flatness is local under localization; (e) in a short exact
sequence, if two terms are flat then so is the third; and (f) a finite module
over a local noetherian ring is flat iff it is free.

\end{proposition}
\begin{proof}

The proof is cited in the text to Atiyah--Macdonald, Chap. 2, and Bourbaki;
the book does not reprove these standard module-theoretic criteria here.

\end{proof}
\begin{definition}[Flat morphism]\label{ha-ch3-def-flat}
\source{hartshorne-algebraic-geometry:III.9}
A morphism is flat when the induced local ring maps are flat.
\end{definition}
\begin{proposition}[Localization and base change of flatness]\label{ha-ch3-prop-flat-preservation}\dcref{III.9.2}
\source{hartshorne-algebraic-geometry:III.9.2}
The properties of flat modules in Proposition~III.9.1A globalize: localization,
base change, and composition preserve flatness.
\uses{ha-ch3-prop-flat-criteria}
\end{proposition}
\begin{proof}

This follows from the corresponding module statements in
Proposition~III.9.1A, applied on affine charts.

\end{proof}
\begin{proposition}[Base change for quasi-coherent sheaves]\label{ha-ch3-prop-flat-base-change}\dcref{III.9.3}
\source{hartshorne-algebraic-geometry:III.9.3}
Let $f:X\to Y$ be a separated morphism of finite type between noetherian
schemes, let $g:Y'\to Y$ be flat, and let $X'=X\times_Y Y'$ with projection
$g':X'\to X$.  Then, for quasi-coherent $\mathcal F$,
\[
g^*R^if_*\mathcal F\simeq R^ig'_*g'^*\mathcal F.
\]

\end{proposition}
\begin{proof}

The proof is the Cech calculation on affine covers: flat base change tensors
the Cech complex with the new base ring, and separatedness keeps all finite
intersections affine.

\end{proof}
\begin{corollary}[Affine flat base change]\label{ha-ch3-cor-flat-affine}\dcref{III.9.4}
\source{hartshorne-algebraic-geometry:III.9.4}
If $Y=\operatorname{Spec}A$ and $y\in Y$ has residue field $k(y)$, then
\[
H^i(X,\mathcal F)\otimes_A k(y)\simeq H^i(X_y,\mathcal F_y)
\]
under the hypotheses of Proposition~III.9.3.
\uses{ha-ch3-prop-flat-base-change}
\end{corollary}
\begin{proof}

Apply Proposition~III.9.3 to the base change $\operatorname{Spec}k(y)\to Y$.

\end{proof}
\begin{proposition}[Flatness and fiber dimensions]\label{ha-ch3-prop-flat-fibers}\dcref{III.9.5}
\source{hartshorne-algebraic-geometry:III.9.5}
Let $f:X\to Y$ be a morphism of finite type over a field $k$, and assume
$f$ is flat.  Then for every $x\in X$, with $y=f(x)$,
\[
\dim_x(X_y)=\dim_xX-\dim_yY.
\]

\end{proposition}
\begin{proof}

This is the local dimension formula for a finite-type flat algebra, as cited
in the text.

\end{proof}
\begin{corollary}[Flat families over regular curves]\label{ha-ch3-cor-flat-curve}\dcref{III.9.6}
\source{hartshorne-algebraic-geometry:III.9.6}
Let $f:X\to Y$ be flat, of finite type over a field, with $Y$ irreducible.
If $Y$ is regular of dimension one, then the following are equivalent: every
associated point of $X$ maps to the generic point of $Y$, and every
irreducible component of every fiber has the expected dimension.
\uses{ha-ch3-prop-flat-fibers}
\end{corollary}
\begin{proof}

Use the non-zero-divisor criterion for a parameter of the regular curve and
the dimension formula of Proposition~III.9.5.

\end{proof}
\begin{proposition}[Flatness and associated points]\label{ha-ch3-prop-flat-specialization}\dcref{III.9.7}
\source{hartshorne-algebraic-geometry:III.9.7}
Let $f:X\to Y$ be a morphism with $Y$ integral and regular of dimension one.
Then $f$ is flat iff every associated point of $X$ maps to the generic point
of $Y$; if $X$ is reduced, this is equivalent to every irreducible component
of $X$ dominating $Y$.

\end{proposition}
\begin{proof}

This is the associated-point flatness criterion, proved in the text by the
local criterion for flatness.

\end{proof}
\begin{proposition}[Flat closure over a curve]\label{ha-ch3-prop-flat-limitations}\dcref{III.9.8}
\source{hartshorne-algebraic-geometry:III.9.8}
Let $Y$ be a regular integral curve, let $P\in Y$ be closed, and let
$X\subset\mathbf P^n_{Y-P}$ be flat over $Y-P$.  The scheme-theoretic closure
$\overline X\subset\mathbf P^n_Y$ is flat over $Y$.

\end{proposition}
\begin{proof}

The proof is the local associated-point argument given in the text: the
closure has no associated point over $P$, so the parameter of the regular
curve is a non-zero-divisor and the local criterion gives flatness.

\end{proof}
\begin{theorem}[Flat families of projective schemes]\label{ha-ch3-thm-flat-family}\dcref{III.9.9}
\source{hartshorne-algebraic-geometry:III.9.9}
For an integral noetherian base $T$ and a flat projective family
$X\subset\mathbf P_T^n$, the Hilbert polynomial of the fibers is constant.

\end{theorem}
\begin{proof}
The graded coordinate module is flat over the base in each large
degree; III.5.2 and the Hilbert polynomial calculation identify its fiber
dimensions, which therefore agree on the integral base.
\end{proof}
\begin{corollary}[Constancy on connected bases]\label{ha-ch3-cor-flat-family}\dcref{III.9.10}
\source{hartshorne-algebraic-geometry:III.9.10}
For connected bases the Hilbert polynomial is constant.

\end{corollary}
\begin{proof}

The coefficients of the fiber Hilbert polynomial are locally constant by the flat-family theorem.
An integer-valued locally constant function is constant on a connected base.

\end{proof}
\begin{theorem}[Normal families over a nonsingular curve]\label{ha-ch3-thm-normal-family}\dcref{III.9.11}
\source{hartshorne-algebraic-geometry:III.9.11}
An algebraic family of normal varieties parametrized by a nonsingular curve is
a flat family of schemes.
\uses{ha-ch3-lem-hironaka}
\end{theorem}
\begin{proof}

Flatness follows from the local flatness criterion.  At a point of a fiber,
let $t$ be a uniformizer on the curve and $\mathfrak p$ the unique minimal
prime over $(t)$.  The generic-point multiplicity hypothesis makes $t$ a
uniformizer in $A_{\mathfrak p}$, while the fiber's normality makes
$A/\mathfrak p$ normal.  Lemma III.9.12 gives $\mathfrak p=tA$, so the
scheme-theoretic fiber is reduced and equals the prescribed variety.

\end{proof}
\begin{lemma}[Hironaka]\label{ha-ch3-lem-hironaka}\dcref{III.9.12}
\source{hartshorne-algebraic-geometry:III.9.12}
For a local noetherian domain, the relevant intersection of principal ideals
in a flat finite-type algebra is controlled by the generic fiber; this gives
the specialization lemma used in the normal-family result.

\end{lemma}
\begin{proof}

Let $t$ be a parameter of the one-dimensional base and $A$ the local algebra.  Flatness makes $t$
 a non-zero-divisor, while normality of the generic and special fibers identifies the contraction of
 $(t)$ with the intersection of the height-one principal ideals above it.  Krull intersection then
 gives $\bigcap_{n\ge1}t^nA=0$ and the claimed control of the special fiber.

\end{proof}

\begin{corollary}[Normal-family genus]\label{ha-ch3-cor-normal-genus}\dcref{III.9.13}
\source{hartshorne-algebraic-geometry:III.9.13}
In a flat family of normal projective varieties the arithmetic genus is
constant.

\end{corollary}
\begin{proof}

The arithmetic genus is the fixed value of the Hilbert polynomial at $0$ (up to the usual sign and
constant term).  Apply the constancy of the Hilbert polynomial in a flat projective family.

\end{proof}

\section{Smooth morphisms}
\label{ha-ch3-sec10}
\begin{definition}[Smooth morphism]\label{ha-ch3-def-smooth}
\source{hartshorne-algebraic-geometry:III.10}
A finite-type morphism is smooth of relative dimension $n$ when it is flat and
its geometric fibers are nonsingular $n$-folds; equivalently it has the formal
lifting property and locally free relative differentials of rank $n$.
\end{definition}
\begin{proposition}[Differentials]\label{ha-ch3-prop-differentials}\dcref{III.10.1}
\source{hartshorne-algebraic-geometry:III.10.1}
For $X\to Y\to Z$ there is a right-exact sequence
$f^*\Omega_{Y/Z}\to\Omega_{X/Z}\to\Omega_{X/Y}\to0$; smoothness is
characterized by local freeness of the last term together with flatness.

\end{proposition}
\begin{proof}

For affine rings $A\to B\to C$, a $C$-derivation of $C$ over $A$ restricts to a $B$-derivation after
quotienting by derivations coming from $B$; the universal property gives the exact sequence.  Localize
to obtain the sheaf statement.  The infinitesimal lifting criterion identifies smoothness with
flatness and locally free relative differentials of the expected rank.

\end{proof}
\begin{theorem}[Jacobian criterion]\label{ha-ch3-thm-jacobian}\dcref{III.10.2}
\source{hartshorne-algebraic-geometry:III.10.2}
For finite-type schemes over a field, smoothness is equivalent to the expected
rank condition on differentials/Jacobian at every point.

\end{theorem}
\begin{proof}
Present the local algebra by generators and relations; the formal
lifting criterion is equivalent to solving the linearized relation equations,
which is exactly the Jacobian rank condition.
\end{proof}
\begin{lemma}[Regular local algebra]\label{ha-ch3-lem-regular-local}\dcref{III.10.3A}
\source{hartshorne-algebraic-geometry:III.10.3A}
The dimension, embedding dimension, and projective dimension relations for a
local homomorphism give the regularity criterion used in the smoothness proof.

\end{lemma}
\begin{proof}

For a local map $(A,\mathfrak m)\to(B,\mathfrak n)$, the conormal exact sequence compares
$\mathfrak n/\mathfrak n^2$ with $\Omega_{B/A}\otimes k(\mathfrak n)$.  The depth and dimension
equalities in a regular local ring make the Jacobian rank equal to the relative codimension; conversely
that rank equality forces the maximal ideal to be generated by the dimension number of parameters.

\end{proof}
\begin{proposition}[Composition and base change for smooth maps]\label{ha-ch3-prop-smooth-composition}\dcref{III.10.4}
\source{hartshorne-algebraic-geometry:III.10.4}
Composites and base changes of smooth morphisms are smooth.
\uses{ha-ch3-prop-differentials}
\end{proposition}
\begin{proof}

Flatness is stable under composition and tensor product.  The differential sequence in
Proposition~\ref{ha-ch3-prop-differentials} shows that the relative differentials of a composite are
locally free of the sum of the ranks, and base change preserves the same presentation and rank.

\end{proof}
\begin{lemma}[Generic separability]\label{ha-ch3-lem-generic-separable}\dcref{III.10.5}
\source{hartshorne-algebraic-geometry:III.10.5}
A dominant finite-type morphism of integral schemes is generically smooth
after removing the inseparable locus; in characteristic zero this holds
without a separability qualification.

\end{lemma}
\begin{proof}

Choose a separating transcendence basis of the extension of function fields.  The module of
Kähler differentials has the expected dimension after inverting one nonzero determinant of the
Jacobian.  Its nonvanishing locus is dense exactly when the extension is separable; in characteristic
zero every finite extension is separable.

\end{proof}
\begin{proposition}[Generic smoothness]\label{ha-ch3-prop-smooth-nonsingular}\dcref{III.10.6}
\source{hartshorne-algebraic-geometry:III.10.6}
A dominant morphism of nonsingular varieties is smooth on a dense open set.

\end{proposition}
\begin{proof}

At the generic point the induced extension of function fields is separable.  By the generic
separability lemma, a Jacobian minor is nonzero there; its nonvanishing locus is a dense open.  The
Jacobian criterion identifies the restriction over this open with a smooth morphism.

\end{proof}
\begin{corollary}[Generic smoothness over a perfect field]\label{ha-ch3-cor-bertini}\dcref{III.10.7}
\source{hartshorne-algebraic-geometry:III.10.7}
Over a perfect field, a finite-type morphism is smooth on a dense open when its
generic extension is separable.

\end{corollary}
\begin{proof}

The generic separability lemma gives a dense open where the differential rank is maximal.  Perfectness
 makes residue-field extensions separable, so geometric regularity agrees with regularity there; the
 Jacobian criterion then gives smoothness.

\end{proof}
\begin{theorem}[Homogeneous-space smoothness]\label{ha-ch3-thm-homogeneous}\dcref{III.10.8}
\source{hartshorne-algebraic-geometry:III.10.8}
A homogeneous space under an algebraic group is nonsingular, and its general
hyperplane sections satisfy the Bertini conclusion.

\end{theorem}
\begin{proof}

Translation by the group identifies all tangent spaces, so it is enough to check smoothness at the
identity coset, where the stabilizer quotient gives the expected tangent dimension.  Apply generic
smoothness to the universal hyperplane section to obtain the Bertini assertion.

\end{proof}
\begin{corollary}[Bertini theorem]\label{ha-ch3-cor-bertini-section}\dcref{III.10.9}
\source{hartshorne-algebraic-geometry:III.10.9}
A general hyperplane section of a nonsingular projective variety is
nonsingular under the stated characteristic hypotheses.

\end{corollary}
\begin{proof}

Apply the generic smoothness theorem to the universal hyperplane family.  The nonsmooth locus has
 proper image in the dual projective space, so a hyperplane outside that image has a nonsingular
 section.

\end{proof}

\section{Formal functions}
\label{ha-ch3-sec11}
\begin{theorem}[Theorem on formal functions]\label{ha-ch3-thm-formal-functions}\dcref{III.11.1}
\source{hartshorne-algebraic-geometry:III.11.1}
Let $f:X\to Y$ be projective noetherian, $y\in Y$, and $\mathcal F$ coherent.
Writing $X_n=X\times_Y\operatorname{Spec}(\mathcal O_{Y,y}/\mathfrak m_y^{n+1})$,
there is a natural isomorphism
\[\widehat{(R^if_*\mathcal F)}_y\cong\varprojlim_nH^i(X_n,\mathcal F|_{X_n}).\]

\end{theorem}
\begin{proof}
Embed $X$ in projective space, resolve $\mathcal F$ by sums of
twists, and apply the finite generation and eventual vanishing theorem to the
resulting complexes.  Completion commutes with their finite cohomology, and
passage to the inverse limit gives the displayed map and its inverse.
\end{proof}
\begin{corollary}[Completion and base change]\label{ha-ch3-cor-formal-functions}\dcref{III.11.2}
\source{hartshorne-algebraic-geometry:III.11.2}
The completed direct image is compatible with completion and base change.

\end{corollary}
\begin{proof}

Reduce the formal-functions isomorphism modulo $\mathfrak m_y^{n+1}$.  The inverse system is then the
 cohomology of the corresponding thickening, which is exactly the base-change complex; passing to the
 inverse limit gives the compatibility.

\end{proof}
\begin{corollary}[Affine formal functions]\label{ha-ch3-cor-formal-affine}\dcref{III.11.3}
\source{hartshorne-algebraic-geometry:III.11.3}
In the affine case, the completed direct-image modules are finite.

\end{corollary}
\begin{proof}

On an affine base, the cohomology modules in the theorem on formal functions are finite by
Serre finiteness.  Their inverse limit is therefore a finitely generated module over the completed
local ring.

\end{proof}
\begin{corollary}[Zariski main theorem]\label{ha-ch3-cor-zariski-main}\dcref{III.11.4}
\source{hartshorne-algebraic-geometry:III.11.4}
A proper birational morphism to a normal scheme has
$f_*\mathcal O_X=\mathcal O_Y$.

\end{corollary}
\begin{proof}

The inclusion $\mathcal O_Y\hookrightarrow f_*\mathcal O_X$ is an isomorphism over the generic point.
For a normal affine open $\operatorname{Spec}A\subset Y$, a section on $f^{-1}\operatorname{Spec}A$
lies in the fraction field and is regular at every codimension-one point because $f$ is proper and
birational.  Normality gives $A=\bigcap_{\operatorname{ht}\mathfrak p=1}A_{\mathfrak p}$, so the
section belongs to $A$.

\end{proof}
\begin{corollary}[Stein factorization]\label{ha-ch3-cor-stein}\dcref{III.11.5}
\source{hartshorne-algebraic-geometry:III.11.5}
A proper morphism factors through a finite morphism after taking the relative
Spec of its direct-image algebra.

\end{corollary}
\begin{proof}

Set $Y'=\operatorname{Spec}_Yf_*\mathcal O_X$.  Properness and coherence make $f_*\mathcal O_X$
 a finite $\mathcal O_Y$-algebra, so $Y'\to Y$ is finite.  The adjunction map $X\to Y'$ is proper
 with geometrically connected fibers, and its composite with $Y'\to Y$ is $f$.

\end{proof}

\section{Semicontinuity}
\label{ha-ch3-sec12}
\begin{proposition}[Universal complexes]\label{ha-ch3-prop-universal-complex}\dcref{III.12.1}
\source{hartshorne-algebraic-geometry:III.12.1}
For a bounded complex of finite free modules over a noetherian ring, its
cohomology after tensoring with a residue field is represented by a bounded
complex of finite free modules with differentials in the maximal ideal.

\end{proposition}
\begin{proof}

Over a local ring, split off every unit entry occurring in a differential.  Each such entry gives a
contractible two-term summand.  Removing all of them leaves a homotopy-equivalent complex whose
 differential matrices have entries in the maximal ideal; the construction localizes and glues.

\end{proof}
\begin{proposition}[Fiber dimensions of complexes]\label{ha-ch3-prop-fiber-dimensions}\dcref{III.12.2}
\source{hartshorne-algebraic-geometry:III.12.2}
The resulting fiber dimensions are controlled by the ranks of the
differentials.

\end{proposition}
\begin{proof}

For a complex of vector spaces, $\dim H^i=\dim C^i-\operatorname{rank}d^{i-1}-\operatorname{rank}d^i$.
Apply this identity after tensoring the universal complex with $k(y)$.

\end{proof}
\begin{lemma}[Minimal complexes]\label{ha-ch3-lem-minimal-complex}\dcref{III.12.3}
\source{hartshorne-algebraic-geometry:III.12.3}
Over a local ring, a complex with differentials in the maximal ideal has fiber
cohomology whose dimension is the rank of the corresponding free term minus
the ranks of adjacent differentials.

\end{lemma}
\begin{proof}

The differentials vanish after tensoring with the residue field, so the fiber cohomology is the
corresponding vector-space kernel modulo image.  Its dimension is the alternating rank expression in
the preceding proposition, and Nakayama lifts the rank conditions back to the local ring.

\end{proof}
\begin{proposition}[Equivalent complex conditions]\label{ha-ch3-prop-complex-conditions}\dcref{III.12.4}
\source{hartshorne-algebraic-geometry:III.12.4}
For a complex of finite free modules, the following are equivalent: a fiber
cohomology dimension is locally constant; the relevant cohomology module is
free and commutes with residue-field base change; and the neighboring
differential has the required constant rank.

\end{proposition}
\begin{proof}

Localize at a point and replace the complex by the minimal complex of III.12.1.  The rank formula of
III.12.2 shows that constant fiber dimension is equivalent to constant rank of the neighboring
differential.  A constant-rank map between finite free modules has locally free kernel and cokernel,
and tensoring with residue fields preserves them; the converse follows by taking dimensions.

\end{proof}
\begin{proposition}[Natural fiber map]\label{ha-ch3-prop-fiber-map}\dcref{III.12.5}
\source{hartshorne-algebraic-geometry:III.12.5}
There is a natural map from the cohomology of a finite complex tensored with
$k(y)$ to the cohomology of the tensor complex, and its kernel/cokernel are
controlled by the neighboring cohomology module.

\end{proposition}
\begin{proof}

Tensor the short exact sequences defining cycles and boundaries with $k(y)$ and use the resulting
Tor exact sequence.  The edge map from $H^i(C^\bullet)\otimes k(y)$ to $H^i(C^\bullet\otimes k(y))$
has kernel the image of $\operatorname{Tor}_1(H^{i+1}(C^\bullet),k(y))$ and cokernel a quotient of
$H^{i+1}(C^\bullet)[\mathfrak m_y]$.

\end{proof}
\begin{corollary}[Local freeness criterion]\label{ha-ch3-cor-local-free-cohomology}\dcref{III.12.6}
\source{hartshorne-algebraic-geometry:III.12.6}
Constancy of the relevant fiber dimension makes the cohomology module locally
free and makes the fiber map an isomorphism.
\uses{ha-ch3-prop-complex-conditions}
\end{corollary}
\begin{proof}

Use Proposition~\ref{ha-ch3-prop-complex-conditions}: constancy gives the required constant rank,
so the cohomology module is locally free and the Tor terms in the natural fiber map vanish.

\end{proof}
\begin{definition}[Upper semicontinuity]\label{ha-ch3-def-upper-semicontinuity}\dcref{III.12.7}
\source{hartshorne-algebraic-geometry:III.12.7}
A function is upper semicontinuous when each locus where its value is at most a
fixed integer is open.
\end{definition}
\begin{theorem}[Semicontinuity]\label{ha-ch3-thm-semicontinuity}\dcref{III.12.8}
\source{hartshorne-algebraic-geometry:III.12.8}
For projective $f:X\to Y$ with $Y$ noetherian and coherent $\mathcal F$ flat
over $Y$, $y\mapsto\dim_{k(y)}H^i(X_y,\mathcal F_y)$ is upper semicontinuous.

\end{theorem}
\begin{proof}
Represent $R^if_*\mathcal F$ locally by a finite complex of free
modules.  The preceding rank description shows that the locus where fiber
cohomology has dimension at most $N$ is open.
\end{proof}
\begin{corollary}[Grauert and base change]\label{ha-ch3-cor-grauert}\dcref{III.12.9}
\source{hartshorne-algebraic-geometry:III.12.9}
If a fiber cohomology dimension is constant near $y$, $R^if_*\mathcal F$ is
locally free there and commutes with arbitrary residue-field base change.

\end{corollary}
\begin{proof}

The semicontinuity theorem makes the constant-dimension locus open.  On that neighborhood apply the
local freeness criterion to the finite free complex representing $R^if_*\mathcal F$.

\end{proof}
\begin{proposition}[Cohomology and base change criterion]\label{ha-ch3-prop-cohomology-base-change}\dcref{III.12.10}
\source{hartshorne-algebraic-geometry:III.12.10}
The base-change map $(R^if_*\mathcal F)\otimes k(y)\to H^i(X_y,\mathcal F_y)$
is surjective precisely under the neighboring-map condition; if it is
surjective, it is an isomorphism near $y$ and $R^{i-1}f_*\mathcal F$ is locally
free there.
\uses{ha-ch3-prop-fiber-map}
\end{proposition}
\begin{proof}

Represent the derived direct image locally by a finite free complex.  The natural fiber map is the
edge map in the Tor exact sequence of Proposition~\ref{ha-ch3-prop-fiber-map}; surjectivity is
equivalent to vanishing of the neighboring Tor term.  Nakayama then makes the neighboring differential
constant rank near $y$, so the map is an isomorphism there and the previous cohomology module is free.

\end{proof}
\begin{theorem}[Cohomology and base change]\label{ha-ch3-thm-cohomology-base-change}\dcref{III.12.11}
\source{hartshorne-algebraic-geometry:III.12.11}
Let $f:X\to Y$ be projective, $Y$ noetherian, and $\mathcal F$ coherent and
flat over $Y$.  The natural map
$(R^if_*\mathcal F)\otimes k(y)\to H^i(X_y,\mathcal F_y)$ is an isomorphism
near $y$ whenever the degree-$i$ map is surjective at $y$; in that case the
preceding direct image is locally free near $y$.

\end{theorem}
\begin{proof}
Apply III.12.10 to the finite free complex representing direct
image cohomology and use the rank/minimal-complex criterion of III.12.4.
\end{proof}
